% =================================================================
\chapter{Technischer Aufbau}
% =================================================================

mLogistics verwendet \textit{Log4j2}, um Logdateien zu erzeugen und zu verwalten.

\section{Logging in Clusterumgebungen}

Log4j2 ist ein leistungsfähiges und flexibles Logging-Framework für Java-Anwendungen und die Weiterentwicklung von Log4j (Apache Logging Services).  
Es ermöglicht Entwicklern, Protokollnachrichten effizient zu erfassen, zu formatieren und zu verwalten.  

Log4j2 ist bekannt für seine hohe Performance, Modularität und einfache Konfiguration.

\subsection{Wichtige Merkmale von Log4j2}

\begin{itemize}
    \item Asynchrones Logging (über LMAX-Disruptor)
    \item Pluggable Architektur
    \item Flexible Konfiguration (YAML, JSON, XML, Java)
    \item Filtermechanismen auf verschiedenen Ebenen
    \item Rollover und Archivierung von Logdateien
    \item Komprimierung
    \item Integration mit diversen Frameworks
\end{itemize}

\subsection{Asynchronität}

Die Protokollierung erfolgt asynchron.  
Das bedeutet, dass das ELS und das Logging-Modul unabhängig voneinander arbeiten, um Performance-Einbußen zu vermeiden.

\section{Root-Verzeichnis}

Standardmäßig werden Logdateien in der im System definierten Verzeichnisstruktur abgelegt.  
Diese Struktur bildet die Basis für konsistente Ablage und Zugriff auf Logs innerhalb der gesamten Umgebung.

\section{Archivierung / Rollover}

% Platzhalter
Hier folgt die detaillierte Beschreibung der Archivierung und der Rollover-Strategien.

\section{Weitere technische Grundlagen}

% Platzhalter
Weitere technische Inhalte und Details zur Implementierung des Loggings werden hier ergänzt.